\section{What the construction guarantees, and what the tails choose}
\label{sec:validity}

The construction separates properties that are often mixed together in a
smile fit.  Positive density is automatic for every finite coefficient
vector.  A finite forward imposes one exact restriction on an exponential
right tail.  The two tail exponents are modeling choices: they decide the
moment domain and the continuation of the smile beyond observable strikes.

\subsection{One-expiry no-arbitrage}

\begin{proposition}[Validity of one slice]
\label{prop:noarb}
Let the parameters satisfy \eqref{eq:martcondition}, and choose the shift of
\cref{prop:shift}.  Then:
\begin{enumerate}[label=(\roman*)]
\item $X$ has a strictly positive continuous density on $\R$;
$Y=e^X$ has a strictly positive density on $(0,\infty)$ and mean one;
\item the call is strictly decreasing and strictly convex in normalized
strike $y$, obeys the bounds in \cref{def:slice}, and satisfies put--call
parity;
\item every nondegenerate finite-width butterfly has positive price;
\item the implied-variance curve satisfies $g_{\rm D}(k)>0$ at every finite
strike.
\end{enumerate}
\end{proposition}

\begin{proof}
Equation \eqref{eq:density} gives a positive continuous density and
\cref{prop:shift} gives unit mean.  In rank coordinates,
\[
c(y)=\int_0^1(e^{Q(u)}-y)^+\,\dd u.
\]
Each integrand is decreasing and convex in $y$.  Differentiation gives
$c'(y)=-\Prob(Y>y)$ and $c''(y)=f_Y(y)>0$.  Jensen's inequality and
dominated convergence give the boundary conditions.  The butterfly and
parity statements follow.  Finally, \eqref{eq:durrleman} expresses the
positive density as a positive Black factor times $g_{\rm D}$.
\end{proof}

This is a statement about the entire continuum of strikes, not merely the
calibration grid.  The numerical representation must preserve the same
geometry; \cref{sec:certificates} gives the corresponding checks.

\begin{figure}[htbp]
\centering
\paperfig{0.92\textwidth}{fig_butterfly.pdf}
\caption{One-slice validity on the fitted SPY December 2026 law.  Left: the
call curve is convex in normalized strike.  Right: symmetric butterfly prices,
divided by squared half-width, converge to the density.  The law supplies the
convexity before calibration begins.}
\label{fig:butterfly}
\end{figure}

\subsection{What the tail exponents mean}
\label{sec:tails}

The log-odds coordinate has a useful probabilistic meaning in the right tail:
$\Prob(Z>z)\sim e^{-z}$.  Hence $z$ is asymptotically the negative logarithm
of tail probability.  If log-return grows linearly with $z$, its tail is
exponential.  If it grows like $\sqrt z$, its tail has Gaussian rate.  The
exponent $\alpha$ moves continuously between these two growth laws.

At the right endpoint, $1-\Lambda(z)\sim e^{-z}$,
$\ell_+(\Lambda(z))=z+1+o(1)$, and $\ell_-(\Lambda(z))=1+o(1)$.
The left endpoint is symmetric.  Integrating \eqref{eq:xspeed} gives
\begin{align}
x(z)&=m+b_+
+\frac{\lambda_+}{1-\alpha_+}(z+1)^{1-\alpha_+}+o(1),
&&z\to+\infty,
\label{eq:righttransporttail}\\
x(z)&=m+b_-
-\frac{\lambda_-}{1-\alpha_-}(1-z)^{1-\alpha_-}+o(1),
&&z\to-\infty,
\label{eq:lefttransporttail}
\end{align}
for finite constants $b_\pm$.  The additive constants do not affect the
leading tail rates.

\begin{proposition}[Log-return tail rates]
\label{prop:tails}
For a generalized LQD slice,
\begin{align}
-\log\Prob(X>x)
&\sim
\left(\frac{(1-\alpha_+)x}{\lambda_+}\right)^{1/(1-\alpha_+)},
&&x\to+\infty,
\label{eq:righttailrate}\\
-\log\Prob(X<-x)
&\sim
\left(\frac{(1-\alpha_-)x}{\lambda_-}\right)^{1/(1-\alpha_-)},
&&x\to+\infty.
\label{eq:lefttailrate}
\end{align}
\end{proposition}

\begin{proof}
Invert \eqref{eq:righttransporttail} and use
$\Prob(Z>z)=1-\Lambda(z)\sim e^{-z}$.  The left proof is identical.
\end{proof}

Write $p=1/(1-\alpha)$.  The three cases are:
\begin{center}
\begin{tabular}{@{}lll@{}}
\toprule
Tail exponent & Log-return tail & Extreme smile wing \\
\midrule
$\alpha=0$ & $-\log\Prob(|X|>x)\asymp x$ & $w(k)\asymp |k|$ \\
$0<\alpha<1/2$ & $-\log\Prob(|X|>x)\asymp x^p$, $1<p<2$
& $w(k)$ grows sublinearly \\
$\alpha=1/2$ & $-\log\Prob(|X|>x)\asymp x^2$ & $w(k)$ tends to a constant \\
\bottomrule
\end{tabular}
\end{center}
At $\alpha=0$, $Y=e^X$ has a power tail.  At $\alpha=1/2$, a normal tail
with standard deviation $s$ corresponds to $\lambda=s/\sqrt2$.
The two sides may use different exponents.  Equity smiles often retain
$\alpha_-=0$ on the crash side and use either $0$ or a small positive value
on the right.  The choice is a tail convention unless far-wing instruments
identify it.

\begin{figure}[htbp]
\centering
\paperfig{\textwidth}{fig_tail_family.pdf}
\caption{The tail exponent.  Left: the leading negative log-tail grows like
$x^{1/(1-\alpha)}$.  Center: the extreme quantile grows like
$z^{1-\alpha}$ in log-odds.  Right: for $\alpha>0$, total implied variance
grows like $|k|^{(1-2\alpha)/(1-\alpha)}$; at $\alpha=1/2$ it approaches a
constant.  Curves are normalized at one to compare rates, not levels.}
\label{fig:tailfamily}
\end{figure}

\subsection{Moments and the finite-forward wall}

Define the two moment boundaries
\begin{equation}
\overline r_\pm=
\begin{cases}
1/\lambda_\pm,&\alpha_\pm=0,\\
+\infty,&\alpha_\pm>0.
\end{cases}
\label{eq:momentboundaries}
\end{equation}

\begin{proposition}[Moment domain]
\label{prop:strip}
For every generalized LQD slice,
\begin{equation}
\E[e^{rX}]<\infty
\quad\Longleftrightarrow\quad
-\overline r_-<r<\overline r_+.
\label{eq:strip}
\end{equation}
Thus the last finite positive moment of $Y$ beyond its mean is
$r_+^*=\overline r_+-1$, and the last finite inverse moment is
$r_-^*=\overline r_-$, with $+\infty$ allowed.
\end{proposition}

\begin{proof}
Write the moment as
\[
\E[e^{rX}]=\int_{-\infty}^{\infty}e^{rx(z)}\rho(z)\,\dd z.
\]
If $\alpha_+>0$, then $x(z)=o(z)$ and the factor $\rho(z)\sim e^{-z}$
dominates every fixed positive $r$.  If $\alpha_+=0$, the right integrand is
asymptotic to a constant times $e^{(r\lambda_+-1)z}$.  The left endpoint
gives the other inequality in the same way.
\end{proof}

The martingale condition is the single point $r=1$ in this interval.  When
$\alpha_+=0$, it is available exactly when $\lambda_+<1$.  A positive right
tail exponent makes every positive moment finite; a positive left exponent
makes every inverse moment finite.  This is why two fits that are virtually
indistinguishable on listed options can behave very differently for power
payoffs, inverse products, and tail-sensitive risk measures.

\subsection{From tail rate to smile wing}

For exponential tails, Lee's moment formula gives the limiting slopes of
total implied variance \citep{Lee2004}:
\begin{equation}
\beta_R=\lim_{k\to+\infty}\frac{w(k)}k,
\qquad
\beta_L=\lim_{k\to-\infty}\frac{w(k)}{|k|}.
\label{eq:lee}
\end{equation}

\begin{proposition}[Exponential-tail wing slopes]
\label{prop:leeclosed}
If $\alpha_+=0$ and $\lambda_+<1$, and if $\alpha_-=0$, then
\begin{equation}
\beta_R=\frac{2\lambda_+}{(1+\sqrt{1-\lambda_+})^2},
\qquad
\beta_L=\frac{2\lambda_-}{(1+\sqrt{1+\lambda_-})^2}.
\label{eq:leeclosed}
\end{equation}
Both maps increase with the corresponding tail scale.  The scale $\lambda$
is not itself the smile slope: \eqref{eq:leeclosed} is the exact conversion.
For small $\lambda$, $\beta\sim\lambda/2$; as $\lambda_+\uparrow1$,
$\beta_R\uparrow2$.
\end{proposition}

\begin{proof}
Lee's function is
$\Psi(r)=2-4(\sqrt{r^2+r}-r)$.  Substitute
$r_+^*=1/\lambda_+-1$ and $r_-^*=1/\lambda_-$ from
\cref{prop:strip}, then rationalize.  The regular variation in
\cref{prop:tails} turns Lee's limsup into a limit
\citep{BenaimFriz2009}.
\end{proof}

Positive tail exponents have zero Lee slope because all moments are finite.
Their sublinear wing growth remains explicit.

\begin{proposition}[Intermediate and Gaussian-rate wings]
\label{prop:sublinearwings}
If $0<\alpha_+\le1/2$, then, as $k\to+\infty$,
\begin{equation}
w(k)\sim
\frac12\left(\frac{\lambda_+}{1-\alpha_+}
\right)^{1/(1-\alpha_+)}
k^{(1-2\alpha_+)/(1-\alpha_+)}.
\label{eq:rightsublinearwing}
\end{equation}
If $0<\alpha_-\le1/2$, the same formula holds on the left after replacing
$k$ by $|k|$ and the plus parameters by the minus parameters.
\end{proposition}

At $\alpha=1/2$, total variance tends to $2\lambda^2$, the variance of the
matching Gaussian tail.  For $0<\alpha<1/2$, it diverges but does so
sublinearly.  The limit $\alpha\downarrow0$ is nonuniform: every positive
$\alpha$ gives zero Lee slope, while $\alpha=0$ gives the positive slope
\eqref{eq:leeclosed}.  No finite strike strip can reliably distinguish zero
from a sufficiently small positive exponent.  Tail-exponent scenarios are
therefore more stable than an unconstrained estimate for every expiry.

\subsection{Expressiveness at a fixed tail class}

The remaining question is whether the polynomial body imposes a material
restriction once the tail class has been chosen.  For fixed $\alpha_\pm$,
define the residual log-speed of a differentiable target quantile $Q^*$ by
rearranging \eqref{eq:skeleton}:
\begin{equation}
g^*(u)=\log Q^{*\prime}(u)+\log u+\log(1-u)
+\alpha_-\log\ell_-(u)+\alpha_+\log\ell_+(u).
\label{eq:targetg}
\end{equation}

\begin{proposition}[Approximation at fixed tail class]
\label{prop:universal}
Suppose $g^*$ extends continuously to $[0,1]$ and the target law has mean one.
If $\alpha_+=0$, also suppose $e^{g^*(1)}<1$.  Then polynomials of the form
\eqref{eq:gpoly} can be chosen so that $g_N\to g^*$ uniformly.  The
corresponding quantiles converge on compact rank intervals, and their call
curves converge uniformly over all strikes.
\end{proposition}

The proof is in \cref{app:proofs}.  The result says that increasing polynomial
order can approximate any smooth law with the selected endpoint behavior;
the tail skeleton, not the polynomial basis, is the substantive asymptotic
restriction.

No finite-order slice has an atom, a zero-density interval, bounded support,
or mass at zero.  Those features require a quantile jump, a flat part, or a
finite endpoint, none of which is available to a positive smooth transport.
Multimodality is allowed: the density rises where the transport slows and
falls where it accelerates.  Increasing basis order permits more such
changes in the body without altering the stated tail class.
